% GENERATED FILE. Edit canonical lessons, then run npm run generate:books.
% canonical-source-hash: fnv1a64:d18f7dc7a0ed3c0a
% canonical-lessons: JA-C59-boushi, JA-C59-kutsushita, JA-C59-megane, JA-C59-sanjuu, JA-C59-sen

\chapter{Hat, Socks, Glasses, Thirty, Thousand}
\label{ch:ja-hat-socks-glasses-thirty-thousand}
\hlchaptermodality{59}

\begin{chapteropening}
\textbf{By the end of this chapter:} \emph{I can say a hat, socks, glasses, thirty and a thousand.}

Complete the last lesson of chapter 59: I can say a hat, socks, glasses, thirty and a thousand.
\end{chapteropening}

\section[\texorpdfstring{\ja{ぼうし} (boushi)}{ぼうし (boushi)}]{\ja{ぼうし} (boushi) — a hat}

\label{lesson:JA-C59-boushi}



\begin{quote}
Before the new one: say the Japanese for clothes, then the Japanese for shoes.
\end{quote}

\subsection*{You'll want to know: \ja{ぼうし}}
\textbf{\ja{ぼうし}} — \emph{boushi} — \textquotedblleft{}a hat\textquotedblright{}.

\begin{cousinweb}
\textbf{\ja{ぼうし}} (\emph{boushi}) is Sino-Japanese (kango): two Chinese characters read with their on'yomi, readings borrowed from Chinese. They are \emph{bō}, \textquotedblleft{}hat\textquotedblright{} (Middle Chinese \emph{mawH}) and \emph{shi}, \textquotedblleft{}child; a noun ending\textquotedblright{} (\emph{tsiX}) (Unihan, KANJIDIC2; Middle Chinese in Baxter's transcription).
\end{cousinweb}

\begin{grammarlens}[title={What you've built}]
One more word for this chapter.
\end{grammarlens}

\subsection*{Your turn}
\begin{itemize}
\raggedright
  \item \emph{Say it:} \emph{boushi}
  \item \emph{Say it:} \emph{boushi}, once more
  \item \emph{Say it:} say \emph{boushi}
  \item \emph{Recall:} say the Japanese for clothes, then the Japanese for shoes, then say \emph{boushi} again
\end{itemize}

\begin{tcolorbox}[breakable,colback=teal!4,colframe=teal!35!black,title={Before you move on}]
What does \ja{ぼうし} mean? (\textquotedblleft{}A hat\textquotedblright{}.) Say it once more.
\end{tcolorbox}
\section[\texorpdfstring{\ja{くつした} (kutsushita)}{くつした (kutsushita)}]{\ja{くつした} (kutsushita) — socks}

\label{lesson:JA-C59-kutsushita}



\begin{quote}
Before the new one: say the Japanese for shoes, then the Japanese for a hat.
\end{quote}

\subsection*{You'll want to know: \ja{くつした}}
\textbf{\ja{くつした}} — \emph{kutsushita} — \textquotedblleft{}socks\textquotedblright{}.

\begin{cousinweb}
\textbf{\ja{くつした}} (\emph{kutsushita}) is a compound of \emph{kutsu}, \textquotedblleft{}shoe\textquotedblright{}, and native \emph{shita}, \textquotedblleft{}under\textquotedblright{}, both kun'yomi of its kanji (JMdict, KANJIDIC2). \emph{Kutsu} is usually counted native, but whether it is inherited or an early loan is not settled.
\end{cousinweb}

\begin{grammarlens}[title={What you've built}]
One more word for this chapter.
\end{grammarlens}

\subsection*{Your turn}
\begin{itemize}
\raggedright
  \item \emph{Say it:} \emph{kutsushita}
  \item \emph{Say it:} \emph{kutsushita}, once more
  \item \emph{Say it:} say \emph{kutsushita}
  \item \emph{Recall:} say the Japanese for shoes, then the Japanese for a hat, then say \emph{kutsushita} again
\end{itemize}

\begin{tcolorbox}[breakable,colback=teal!4,colframe=teal!35!black,title={Before you move on}]
What does \ja{くつした} mean? (\textquotedblleft{}Socks\textquotedblright{}.) Say it once more.
\end{tcolorbox}
\section[\texorpdfstring{\ja{めがね} (megane)}{めがね (megane)}]{\ja{めがね} (megane) — glasses}

\label{lesson:JA-C59-megane}



\begin{quote}
Before the new one: say the Japanese for a hat, then the Japanese for socks.
\end{quote}

\subsection*{You'll want to know: \ja{めがね}}
\textbf{\ja{めがね}} — \emph{megane} — \textquotedblleft{}glasses\textquotedblright{}.

\begin{cousinweb}
\textbf{\ja{めがね}} (\emph{megane}) is native, inherited rather than borrowed, usually taken apart as \emph{me}, \textquotedblleft{}eye\textquotedblright{}, and \emph{kane}, \textquotedblleft{}metal\textquotedblright{}, voiced to \emph{gane}. JMdict writes it with characters for \textquotedblleft{}eye\textquotedblright{} and \textquotedblleft{}mirror\textquotedblright{}, read as a whole, and marks it as usually written in kana.
\end{cousinweb}

\begin{grammarlens}[title={What you've built}]
One more word for this chapter.
\end{grammarlens}

\subsection*{Your turn}
\begin{itemize}
\raggedright
  \item \emph{Say it:} \emph{megane}
  \item \emph{Say it:} \emph{megane}, once more
  \item \emph{Say it:} say \emph{megane}
  \item \emph{Recall:} say the Japanese for a hat, then the Japanese for socks, then say \emph{megane} again
\end{itemize}

\begin{tcolorbox}[breakable,colback=teal!4,colframe=teal!35!black,title={Before you move on}]
What does \ja{めがね} mean? (\textquotedblleft{}Glasses\textquotedblright{}.) Say it once more.
\end{tcolorbox}
\section[\texorpdfstring{\ja{さんじゅう} (sanjuu)}{さんじゅう (sanjuu)}]{\ja{さんじゅう} (sanjuu) — thirty}

\label{lesson:JA-C59-sanjuu}



\begin{quote}
Before the new one: say the Japanese for socks, then the Japanese for glasses.
\end{quote}

\subsection*{You'll want to know: \ja{さんじゅう}}
\textbf{\ja{さんじゅう}} — \emph{sanjuu} — \textquotedblleft{}thirty\textquotedblright{}.

\begin{cousinweb}
\textbf{\ja{さんじゅう}} (\emph{sanjū}) is Sino-Japanese: \emph{san}, \textquotedblleft{}three\textquotedblright{} (Middle Chinese \emph{sam}), and \emph{jū}, \textquotedblleft{}ten\textquotedblright{} (Middle Chinese \emph{dzyip}), both on'yomi borrowed from Chinese.
\end{cousinweb}

\begin{grammarlens}[title={What you've built}]
One more word for this chapter.
\end{grammarlens}

\subsection*{Your turn}
\begin{itemize}
\raggedright
  \item \emph{Say it:} \emph{sanjuu}
  \item \emph{Say it:} \emph{sanjuu}, once more
  \item \emph{Say it:} say \emph{sanjuu}
  \item \emph{Recall:} say the Japanese for socks, then the Japanese for glasses, then say \emph{sanjuu} again
\end{itemize}

\begin{tcolorbox}[breakable,colback=teal!4,colframe=teal!35!black,title={Before you move on}]
What does \ja{さんじゅう} mean? (\textquotedblleft{}Thirty\textquotedblright{}.) Say it once more.
\end{tcolorbox}
\section[\texorpdfstring{\ja{せん} (sen)}{せん (sen)}]{\ja{せん} (sen) — a thousand}

\label{lesson:JA-C59-sen}



\begin{quote}
Before the new one: say the Japanese for glasses, then the Japanese for thirty.
\end{quote}

\subsection*{You'll want to know: \ja{せん}}
\textbf{\ja{せん}} — \emph{sen} — \textquotedblleft{}a thousand\textquotedblright{}.

\begin{cousinweb}
\textbf{\ja{せん}} (\emph{sen}) is Sino-Japanese (kango): the on'yomi of a character meaning \textquotedblleft{}thousand\textquotedblright{}, a reading borrowed from Chinese (Middle Chinese \emph{tshen} in Baxter's transcription; Unihan, KANJIDIC2).
\end{cousinweb}

\begin{grammarlens}[title={What you've built}]
One more word for this chapter.
\end{grammarlens}

\subsection*{Your turn}
\begin{itemize}
\raggedright
  \item \emph{Say it:} \emph{sen}
  \item \emph{Say it:} \emph{sen}, once more
  \item \emph{Say it:} say \emph{sen}
  \item \emph{Recall:} say the Japanese for glasses, then the Japanese for thirty, then say \emph{sen} again
\end{itemize}

\begin{tcolorbox}[breakable,colback=teal!4,colframe=teal!35!black,title={Before you move on}]
What does \ja{せん} mean? (\textquotedblleft{}A thousand\textquotedblright{}.) Say it once more.
\end{tcolorbox}
